\section*{Chapitre \ref{ch:b1:organomagnesium} --- Les organomagnésiens}

\begin{solution}{exo:b1:organomagnesium:1}
\ce{CH3CH2Br + Mg -> CH3CH2MgBr}~; \ce{C6H5Br + Mg -> C6H5MgBr} (dans
l'éther anhydre). C--Br~: Br (2.96) plus électronégatif que C (2.55),
carbone positif. C--Mg~: Mg (1.31) moins électronégatif, carbone négatif.
\end{solution}

\begin{solution}{exo:b1:organomagnesium:2}
Eau~: \ce{CH3MgI + H2O -> CH4 + Mg(OH)I}.
Éthanol~: \ce{CH3MgI + CH3CH2OH -> CH4 + CH3CH2OMgI}.
Acide éthanoïque~: \ce{CH3MgI + CH3COOH -> CH4 + CH3COOMgI}.
Oxyde de deutérium~: \ce{CH3MgI + D2O -> CH3D + Mg(OD)I}.
\end{solution}

\begin{solution}{exo:b1:organomagnesium:3}
Méthanal~: propan-1-ol (primaire). Éthanal~: butan-2-ol (secondaire).
Propanone~: 2-méthylbutan-2-ol (tertiaire). Benzaldéhyde~:
1-phénylpropan-1-ol (secondaire). Dioxyde de carbone~: acide
propanoïque.
\end{solution}

\begin{solution}{exo:b1:organomagnesium:4}
Son groupe OH est assez acide pour détruire le réactif dès sa
formation~: la molécule se protonerait elle-même (ou ses voisines), en
donnant l'alcoolate de magnésium du butan-1-ol, \ce{CH3(CH2)3OMgBr}, qui
n'est pas un réactif utilisable. Il faut d'abord protéger le groupe OH
(le \cref{ch:b1:acetals-protection} montre l'idée pour un aldéhyde).
\end{solution}

\begin{solution}{exo:b1:organomagnesium:5}
Le doublet de C--Mg attaque le carbone du carbonyle de l'éthanal, le
doublet $\pi$ de C=O passant sur O~: \ce{CH3CH(CH3)O^-} avec
\ce{MgBr+}~; puis \ce{H3O+} cède un proton à l'oxygène~: propan-2-ol. Le
propan-2-ol porte deux groupes méthyle sur le carbone fonctionnel~: il
est \omterm{def:b1:stereochemistry-in-depth:stereocentre}{achiral}. (Avec un autre R, par exemple le bromure d'éthylmagnésium
sur l'éthanal, le produit, le butan-2-ol, est \omterm{def:b1:stereochemistry-in-depth:stereocentre}{chiral} mais racémique~: le
carbonyle plan est attaqué également sur ses deux faces.)
\end{solution}

\begin{solution}{exo:b1:organomagnesium:6}
Le carbone fonctionnel porte un groupe méthyle et deux groupes éthyle~:
le bromure d'éthylmagnésium avec la butan-2-one, ou le bromure de
méthylmagnésium avec la pentan-3-one.
\end{solution}

\begin{solution}{exo:b1:organomagnesium:7}
Préparer \ce{(CH3)2CHMgBr} à partir du 2-bromopropane et de magnésium
dans l'éther anhydre, le verser sur du \ce{CO2} solide, puis acidifier~:
\ce{(CH3)2CHCOOH}. $n = 12.3/123.0 = 0.100$ mol~;
$0.100 \times 0.70 \times 88.0 = \qty{6.2}{g}$.
\end{solution}

\begin{solution}{exo:b1:organomagnesium:8}
$0.18/18.0 = \qty{0.010}{mol}$ d'eau détruisent \qty{0.010}{mol} de
réactif~: \qty{20}{\percent}~; il se forme du benzène,
\ce{C6H5MgBr + H2O -> C6H6 +
Mg(OH)Br}.
\end{solution}

\begin{solution}{exo:b1:organomagnesium:9}
Le réactif, \omterm{def:b1:electronic-effects:nucleophile}{nucléophile}, réagit avec le bromobenzène encore présent~:
bilan \ce{C6H5MgBr + C6H5Br -> C6H5-C6H5 + MgBr2}. Ajouter lentement le
bromure sur un excès de magnésium maintient sa concentration basse~: il
rencontre du métal frais plutôt que le réactif déjà formé.
\end{solution}

\begin{solution}{exo:b1:organomagnesium:10}
1-Phényléthan-1-ol~: \ce{CH3MgBr} (issu du bromométhane) avec le
benzaldéhyde, ou \ce{C6H5MgBr} (issu du bromobenzène) avec l'éthanal.
2-Phénylpropan-2-ol~: \ce{C6H5MgBr} avec la propanone, ou \ce{CH3MgBr}
avec la 1-phényléthan-1-one.
\end{solution}

\begin{solution}{exo:b1:organomagnesium:11}
On ajoute la cétone au réactif afin que celui-ci, préparé dans le ballon
sec, n'en sorte jamais, et que la chaleur dégagée par chaque portion
soit absorbée par l'éther en ébullition. L'eau seule donnerait le sel
basique gélatineux \ce{Mg(OH)Br}, qui piège le produit~; l'acide dilué
froid le dissout sous forme de \ce{Mg^2+}, et le froid limite les
réactions secondaires de l'alcool tertiaire en milieu acide.
\end{solution}

\begin{solution}{exo:b1:organomagnesium:12}
Chaque mole de réactif donne avec l'eau une mole de sel basique de
magnésium, \ce{RMgBr + H2O -> RH + Mg(OH)Br}, dont l'ion \ce{OH-} est
titré par l'acide~: la quantité de réactif est égale à la quantité
d'acide utilisée, \qty{0.088}{mol}, soit un \omterm{def:b1:extent-q-and-k:fractional-extent}{rendement} de
\qty{88}{\percent}.
\end{solution}

\begin{solution}{pb:b1:organomagnesium:1}
\textbf{1.} \ce{C6H5Br + Mg -> C6H5MgBr}.
\textbf{2.} Mg $0.535/24.3 = \qty{0.0220}{mol}$~; \ce{C6H5Br} $3.14/156.9 =
\qty{0.0200}{mol}$~: le magnésium est en léger excès.
\textbf{3.} L'eau détruit le réactif~: \ce{C6H5MgBr + H2O -> C6H6 +
Mg(OH)Br}.
\textbf{4.} Elle empêche la vapeur d'eau de l'air d'entrer par le
réfrigérant.
\textbf{5.} L'éther cède des \omterm{def:b1:lewis-resonance-vsepr:lewis}{doublets non liants} au magnésium (\omterm{def:b1:lewis-resonance-vsepr:lewis-acid}{base de
Lewis})~: sans lui, le réactif ne se forme pas ou ne reste pas en
solution.
\textbf{6.} \ce{C6H5MgBr + C6H5Br -> C6H5-C6H5 + MgBr2}~; biphényle
$0.15/154.0 = \qty{9.7e-4}{mol}$, qui a consommé \qty{9.7e-4}{mol} de
bromobenzène et autant de réactif~: \qty{1.9e-3}{mol} de bromobenzène
perdues au total.
\textbf{7.} $3.28/182.0 = \qty{0.0180}{mol}$.
\textbf{8.} Le doublet de C--Mg attaque le carbone du carbonyle, le
doublet $\pi$ passe sur l'oxygène~: \ce{(C6H5)3C-O^-} \ce{MgBr+}.
\textbf{9.} Réactif disponible au plus $0.0200 - 0.0019 = \qty{0.0181}{mol}$~;
benzophénone \qty{0.0180}{mol}~: la benzophénone est limitante, de peu.
\textbf{10.} Non~: le carbone fonctionnel porte trois groupes phényle
identiques.
\textbf{11.} Le réactif reste dans son ballon sec, et la chaleur de
chaque portion de cétone est absorbée par l'éther en ébullition.
\textbf{12.} Elle détruirait une partie du réactif (benzène) et
laisserait de la benzophénone inutilisée.
\textbf{13.} \ce{(C6H5)3COMgBr + H3O+ -> (C6H5)3COH + Mg^2+ + Br- + H2O}.
\textbf{14.} L'acide dissout les sels basiques de magnésium que l'eau
seule ferait précipiter~; le froid limite la formation du \omterm{def:b1:electronic-effects:intermediates}{carbocation}
de la question 18.
\textbf{15.} Le triphénylméthanol et le biphényle dans la \omterm{def:b1:extent-q-and-k:system}{phase}
éthérée~; les sels de magnésium dans la \omterm{def:b1:extent-q-and-k:system}{phase} aqueuse.
\textbf{16.} En lavant (ou en triturant) le solide avec de l'éther de
pétrole froid~: le biphényle se dissout, le triphénylméthanol reste.
\textbf{17.} Une température de fusion nette, égale à la valeur tabulée
(ou une seule tache en chromatographie sur couche mince).
\textbf{18.} \ce{(C6H5)3COH + H+ -> (C6H5)3C+ + H2O}~: un \omterm{def:b1:electronic-effects:intermediates}{carbocation}
tertiaire dont la charge est étalée par mésomérie sur trois cycles
benzéniques, très stable, et coloré du fait de cette conjugaison étendue.
\textbf{19.} $3.75/260.0 = \qty{0.0144}{mol}$.
\textbf{20.} $0.0144/0.0180 = \qty{80}{\percent}$.
\textbf{21.} Pertes lors des transvasements et de la recristallisation~;
réactif détruit par des traces d'humidité~; réaction incomplète.
\textbf{22.} L'acide benzoïque, \ce{C6H5COOH}, après acidification.
\textbf{23.} $0.0181 \times 122.0 = \qty{2.21}{g}$.
\textbf{24.} \ce{C6H5MgBr + CO2 -> C6H5COOMgBr}, puis
\ce{C6H5COOMgBr + H3O+ -> C6H5COOH + Mg^2+ + Br- + H2O}.
\textbf{25.} \textbf{\qty{80}{\percent}}.
\end{solution}
